%% ch07 --- Imports and dependency injection: modules, cycles and the
%% composition root.
%%
%% Written. The Polish twin must keep the same structure, the same numbers in
%% the same order and the same listing, figure and exercise keys; parity's C4,
%% C12 and C14 compare all three.

\chapter{Imports and dependency injection: modules, cycles and the composition root}
\label{ch:imports}

\noindent
A C\# \code{using} directive is a compile-time instruction about where to
look for names. It runs nothing. Nothing in the assembly it names happens
because you wrote it, the order of the directives at the top of a file does
not matter, and writing one twice is a style question rather than a
behaviour.

None of that is true in Python, and almost everything that goes wrong with
imports is a consequence of the single sentence that replaces it:
\textbf{a module is an object, built by running a file, once}.
\code{import} is the expression that builds it. This chapter takes that
sentence seriously for five sections, and by the end it turns out to have
settled a question that looks unrelated: where dependency injection goes
when the language has no container.

\begin{csbox}
The closest C\# construction is not \code{using} at all. It is a static
class with a static constructor: one instance for the process, initialised
by running code, on first contact. Python's version differs in two ways
that matter. It is triggered by \emph{import} rather than by first use, and
the language makes no thread-safety promise about it that you should rely
on.
\end{csbox}

\section{A module is an object that runs once}
\label{sec:module-object}
\index{modules!as objects}

Here is a module whose body does something you can see.

\pyfile{code/ch07/greeting.py}{A module body is code, and it runs}{lst:ch07-greeting}

\noindent
Now a second file imports it three times. Before you run it: how many times
does the body print, and are the three names the same object?

\pyfile{code/ch07/runs_once.py}{Three imports of one module}{lst:ch07-runs-once}

\transcript{ch07-runs-once}

\noindent
Once, and yes. The mechanism is a dictionary. \code{import x} looks for
\code{"x"} in \api{sys.modules}; a miss runs the file top to bottom and puts
the resulting module object in the cache under that key, and a hit returns
what is already there and runs nothing at all.

\api{importlib.import\_module} is the same operation spelled as a function,
which is why the listing above can use it to stand in for the second and
third importers.

\mermaidfig{ch07-import-once}{What \code{import x} does. The body of \code{x} runs on the miss, and never again}{the import statement as a cache lookup}

\noindent
Three consequences follow immediately, and all three are things engineers
arriving from .NET get wrong.

\begin{itemize}[leftmargin=1.4em, itemsep=3pt]
  \item \textbf{Importing is not free.} Whatever the body does \dash{}
        reading a file, opening a connection, asking the environment a
        question \dash{} it does on somebody's first import, which is
        usually the slowest command you own rather than the one that needed
        it.
  \item \textbf{A module-level name is process-global.} There is one object,
        so a mutable bound at module level is shared by every importer. It
        is a static field, with none of the ceremony that makes a static
        field look like a decision.
  \item \textbf{Import order becomes observable.} \code{SALUTATION} in
        Listing~\ref{lst:ch07-greeting} reads the environment at import
        time, so what it holds depends on when the first import happened,
        and a test that sets the variable afterwards changes nothing.
\end{itemize}

\begin{trapbox}
\textbf{\enquote{A module is a namespace; importing it is free and pure.}}
It is an object, and building it runs a file. A side effect at module level
runs once for the whole process, at a moment decided by import order rather
than by you \dash{} and the cheapest fix is almost always to move the work
into a function, where the caller decides when it happens.
\end{trapbox}

\begin{exercise}{e07_01_import_time_work}{Move the work out of import time}
The starter builds a settings dictionary from the environment at module
level. Delete it and give the module a \code{build\_settings()} function
that does the same work when it is called. One of the tests imports the
module, changes the environment and \emph{then} calls the function: a
module-level dict comprehension cannot pass it, and that is the whole point.
\end{exercise}

\section{Packages, and the two ways to run a file}
\label{sec:packages}
\index{modules!packages}

A package is a directory with an \code{\_\_init\_\_.py}, and the useful way
to think about it is that a package is a module too: \code{\_\_init\_\_.py}
\emph{is} the package's body, it runs once before anything inside the
package, and whatever it binds is what \code{from shop import x} can find.
An \code{\_\_init\_\_.py} that imports its own submodules for the
convenience of callers is charging every importer of any one submodule for
all of them, which is why the default in this book is that the file is
empty.

That leaves the two questions the reader has actually met: absolute against
relative imports, and why \code{python -m} exists. They have one answer.
This module reports what the interpreter thinks it is.

\pyregion{code/ch07/shop/where.py}{relative}{A relative import, attempted at import time}{lst:ch07-relative}

\noindent
And this one runs the same file both ways.

\pyfile{code/ch07/two_ways.py}{The same file, run as a file and as a module}{lst:ch07-two-ways}

\transcript{ch07-two-ways}

\noindent
Run as a file, the module has no package \dash{} \code{\_\_package\_\_} is
\code{None} \dash{} so a relative import has nothing to be relative
\emph{to}, and the failure is an \code{ImportError} about a parent package
rather than a missing file. Run with \code{-m}, the interpreter imports it
as part of \code{shop} and the relative import resolves. Note the other
difference, because it is the one that bites later: the head of
\api{sys.path} is the \emph{script's own directory} in the first case and
the directory you ran from in the second.

\begin{csbox}
\code{-m} is closest to running a project rather than an assembly file:
\code{dotnet run} resolves the project and its references, where invoking a
DLL directly gives you whatever the directory happens to contain. The
\code{src/} layout of Chapter~\ref{ch:packaging} exists to make the second
kind of accident impossible, and this section is the mechanism underneath
it.
\end{csbox}

\noindent
That \api{sys.path} entry comes \emph{first}, so a file of yours wins over
anything with the same name further down, the standard library included.
The four lines below are not wrong. Copy them into a directory under one
name and they work; under another name they do not.

\pyfile{code/ch07/shadow/innocent.py}{Four lines that do not depend on their file name}{lst:ch07-innocent}

\pyfile{code/ch07/shadowing.py}{The same four lines, under two names}{lst:ch07-shadowing}

\transcript{ch07-shadowing}

\noindent
\code{import random} inside a file called \code{random.py} finds the file
itself, binds the module to the name, and the module has no \code{random}
function on it. The error names neither the shadowing nor the file that
caused it.

\begin{trapbox}
\textbf{\enquote{I can name a file, or a variable, \code{list}, \code{id},
\code{type} or \code{random}.}} At module level the name wins over the
standard library for everything below it on \api{sys.path}; inside a
function it wins over the builtin for the rest of the scope. Both fail some
distance from the line that caused them, and the message names the symptom.
\end{trapbox}

\noindent
The last piece of a module's public surface is the one C\# spells with an
access modifier. Python has no \code{internal}: a leading underscore is a
convention, and \code{\_\_all\_\_} is the only thing that makes
\code{from module import *} mean anything definite. Without it the star
copies every name not starting with an underscore, which includes the
modules that file imported \dash{} so a star import can hand you somebody
else's \code{os}.

\begin{exercise}{e07_02_public_surface}{Say what the module exports}
Give the starter an \code{\_\_all\_\_} naming only the two functions
callers outside it should use. The tests check what a star import would
copy, and also that the helpers are still reachable by name, because
\code{\_\_all\_\_} governs the star and hides nothing.
\end{exercise}

\begin{trapbox}
\textbf{\enquote{\code{from x import *} is \code{using x;}.}} It copies
names into your module rather than making them visible, so where a name came
from stops being answerable by reading the file. Prefer \code{import x}, or
\code{from x import name} for a handful of names you use constantly.
\end{trapbox}

\section{A circular import is not the error you think}
\label{sec:cycles}
\index{modules!circular imports}
\index{troubleshooting!circular import}

Two modules that import each other is the shape everybody has been told to
fear. Before reading on, decide which of these fails: a pair where each does
\code{import} of the other, or a pair where each does \code{from} the other
\code{import} a name.

The first runs. The cycle is never the error by itself, because the module
\emph{object} is put in \api{sys.modules} before its body runs \dash{} that
is the order, and it is what stops the recursion. What fails is asking a
half-run module for something it has not bound yet, and \code{from x import
y} asks for exactly that, at import time.

\pyfile{code/ch07/cycles.py}{The same cycle, four ways}{lst:ch07-cycles}

\transcript{ch07-cycles}

\mermaidfig{ch07-cycle-break}{Where a cycle breaks, and where it does not: a module object binds at once, a name in it does not}{the cycle, and the line at which it fails}

\noindent
So the three ways out are all ways of moving the lookup past import time,
and only the last of them removes the cycle.

\begin{enumerate}[leftmargin=1.6em, itemsep=3pt]
  \item \textbf{Hold the module, not the name.} \code{import x}, then
        \code{x.NAME} where it is used. Binding the module always succeeds;
        by the time anything reads an attribute off it, the body has
        finished.
  \item \textbf{Import inside the function.} After the first call it is a
        dictionary lookup, and it takes the whole question out of import
        order. The cost is that a reader has to go and find it, so put it in
        the smallest scope that works.
  \item \textbf{Extract what both sides wanted into a third module.} This is
        the one that is usually right, and a cycle is usually evidence that
        it was owed: two modules that need each other generally share a
        third thing neither of them owns.
\end{enumerate}

\begin{warning}
\textbf{The diagnosis depends on where the file sits, not on what went
wrong.} The last two entries in the transcript above are the same fault.
Inside a package, CPython says \emph{partially initialized module
\dots{} most likely due to a circular import}, which names the cause. In
two files beside the script you ran, it says \emph{consider renaming
\dots{} if it has the same name as a library you intended to import},
which is about \S\ref{sec:packages}'s trap and not this one. Both were run
to produce that transcript. The second message fires exactly in the layout a
first Python project takes, so if you have ever chased a name clash that did
not exist, this is why.
\end{warning}

\begin{trapbox}
\textbf{\enquote{A circular import is an error.}} It is a fact about your
module graph, and Python tolerates it. The error is a \emph{name} read from
a module whose body has not reached it, which is why changing
\code{from x import y} to \code{import x} sometimes fixes a build with no
other change \dash{} and why that fix, taken alone, leaves the cycle there.
\end{trapbox}

\begin{exercise}{e07_03_defer_the_import}{Pay for an import when it is used}
The starter imports a standard-library module at the top and uses it in one
function most callers never reach. Move the import inside that function. The
tests clear the cache, import the module and check that nothing arrived,
then call the function and check that it did.
\end{exercise}

\section{The composition root is a function}
\label{sec:composition-root}
\index{dependency injection!composition root}

Nothing in the previous three sections is about dependency injection, and
all of it is. A container exists to stop objects reaching out for their
collaborators; in Python the thing they reach for is a module-level name,
which \S\ref{sec:module-object} showed is process-global and bound at an
awkward time. Take that away and what is left needs no machinery.

Start with what the work needs, expressed as a \code{Protocol}: structural,
so a class satisfies it by shape and never says so.

\pyregion{code/ch07/composition.py}{ports}{Two protocols: what the service needs, and nothing more}{lst:ch07-ports}

\noindent
The work takes its collaborators as arguments, like anything else.

\pyregion{code/ch07/composition.py}{service}{The service names no concrete type}{lst:ch07-service}

\noindent
The implementations are ordinary classes. Neither inherits from the
protocol, and neither mentions it.

\pyregion{code/ch07/composition.py}{adapters}{The real implementations, which declare nothing}{lst:ch07-adapters}

\noindent
And the composition root is one function: the only place in the program
where a concrete class is named.

\pyregion{code/ch07/composition.py}{root}{The composition root, entire}{lst:ch07-root}

\mermaidfig{ch07-composition-root}{Wiring without a container: one function at the edge knows the concrete types, and nothing else does}{the composition root}

\noindent
\api{functools.partial} binds the leading arguments and leaves the rest of
the signature alone, so what comes back is the service with its wiring
already done. That is what a container hands you, assembled by hand, in one
line you can read \dash{} and because \code{build()} is annotated as
returning a plain callable, nothing downstream can reach past it to ask what
it really is. A checker enforces that.

\begin{csbox}
\code{IServiceCollection} is a registry, and the registration is a second
description of the program that has to be kept true. A composition root is
the program: if a collaborator is missing, the call does not compile under a
checker and does not run under a test. Nothing is resolved at run time, so
there is no such thing here as a container that starts and then throws on
the first request.
\end{csbox}

\begin{exercise}{e07_04_composition_root}{Write the composition root}
The starter has the protocols, the work and a \code{build} that raises.
Write it. The test then wires the service with a fake clock and a list,
neither of which inherits from anything, and checks that two separate
wirings share nothing.
\end{exercise}

\section{Lifetimes, and the container you will actually meet}
\label{sec:lifetimes}
\index{dependency injection!lifetimes}

Three lifetimes are what a .NET engineer will miss first. Each has a Python
answer, and none of them is a registration.

\begin{center}
\small
\begin{tabularx}{\linewidth}{@{}lX@{}}
\toprule
\textbf{C\#} & \textbf{Python} \\
\midrule
\code{AddSingleton} & An object built once in the composition root and
  passed down. A module-level object is the other way to do it and is the
  one to be careful with: it is built at import time, by
  \S\ref{sec:module-object}'s rules. \\
\code{AddScoped} & A parameter. The scope is whatever function holds the
  object \dash{} a request handler, a task, a test \dash{} and a
  \code{with} block closes it. \\
\code{AddTransient} & Call the factory. \code{partial} binds a factory as
  readily as it binds an instance. \\
\code{IServiceProvider} & Usually nothing. Where a framework needs one, it
  brings its own. \\
\bottomrule
\end{tabularx}
\end{center}

\noindent
Which leaves the honest question the brief for this chapter asked: is there
a container worth using? Searched on PyPI rather than remembered, in
\pinnedon{}, the answer is that several are maintained and actively
released \dash{} \pkg{dependency-injector}, \pkg{injector}, \pkg{punq},
\pkg{svcs}, \pkg{wireup} and \pkg{kink} among them \dash{} and that none of
them is the one everybody uses. That is the finding. In .NET the container
ships with the platform, so using it is the default and not a decision; in
Python, adopting one is a decision you will have to defend to the next
person, against a composition root that is a function.

\begin{trapbox}
\textbf{\enquote{I need a DI container.}} You need the thing a container
was bought for, which is that objects do not reach out for their
collaborators. A function and \api{functools.partial} do that, a
\code{Protocol} types it, and a test substitutes anything without
registering it. Reach for a container when you can name what it would buy
you here \dash{} and note that the one most readers will actually meet is
not a general-purpose container at all.
\end{trapbox}

\noindent
That one is FastAPI's \api{Depends}, and it is a container with a very
small remit: it wires the arguments of a request handler, per request. It is
worth previewing here only to say what it is, because everything that makes
it interesting \dash{} what its lifetimes really are, what a \code{yield}
dependency does, where application-lifetime state goes \dash{} belongs to
Chapter~\ref{ch:services}, which is where the framework is.

\begin{versionbox}
\api{Depends} in \pkg{fastapi}~\fastapiver{} takes the callable plus
\code{use\_cache} and \code{scope}. Read that signature out of the installed
package rather than out of a tutorial, because a tutorial is a claim about
some version; Chapter~\ref{ch:services} does.
\end{versionbox}

\noindent
The payoff the chapter promised is now available to state in one sentence,
and it is a structural property rather than a style: \textbf{a package in
which no module mutates anything global at import time, and in which the
only function naming a concrete class is the one at the edge, is a package
in which a test can substitute anything without patching}. Both halves are
this chapter's. The first is what \S\ref{sec:module-object} to
\S\ref{sec:cycles} were for; the second is \S\ref{sec:composition-root},
which is thirty lines long because the first half did the work.
